% Sections/06-HienThucVaTrienKhaiHeThong.tex
% Chương 6 -- Hiện thực và triển khai hệ thống
%=========================

\section{Hiện thực và triển khai hệ thống}

Chương này trình bày kết quả hiện thực và cách triển khai hệ thống thương mại điện tử tích hợp trợ lý mua sắm đa phương thức. Nội dung kết nối phần hiện thực giao diện, dữ liệu và các thành phần AI với quá trình đóng gói, cấu hình, triển khai và vận hành. Các quyết định kiến trúc và nghiệp vụ đã được trình bày ở Chương~\ref{sec:system-design}; chương này tập trung vào cách các quyết định đó được hiện thực trong môi trường chạy thực tế.

% \subsection{Tổng quan và phạm vi hiện thực}

% Hệ thống được hiện thực theo các nhóm thành phần gồm Frontend React, Spring Boot Backend, AI Service FastAPI, MySQL, Qdrant và các dịch vụ hạ tầng hỗ trợ. Frontend cung cấp các màn hình mua sắm, tài khoản, trợ lý AI và quản trị; Backend điều phối xác thực, nghiệp vụ và dữ liệu; AI Service thực hiện embedding, truy hồi, reranking và hội thoại. MySQL giữ dữ liệu nghiệp vụ chuẩn, còn Qdrant lưu chỉ mục dẫn xuất phục vụ tìm kiếm.

% Quy trình hiện thực gồm chuẩn bị dữ liệu, xây dựng các image dịch vụ, cấu hình biến môi trường, chạy kiểm tra ở local và triển khai bản pilot trên Google Cloud Platform. Mỗi thành phần được đóng gói độc lập để có thể cập nhật hoặc rollback theo image tag; dữ liệu lâu dài được tách khỏi vòng đời container bằng named volume hoặc cơ chế lưu trữ tương ứng.

\subsection{Hiện thực giao diện và các luồng nghiệp vụ}

Phần hiện thực giao diện tập trung vào các luồng mà khách hàng và quản trị viên trực tiếp sử dụng. Các màn hình, cách liên kết giữa chúng và nguyên tắc phản hồi được trình bày theo nhóm ở nội dung chính của chương.

% Fragment giao diện được input trong
% Sections/06-HienThucVaTrienKhaiHeThong.tex -- Hiện thực và triển khai hệ thống.

% Các ảnh giao diện được trình bày theo nhóm để giữ nội dung trong chương chính.
\newcommand{\frontendimage}[2]{%
    \IfFileExists{#1}{%
        \includegraphics[width=\linewidth,height=#2,keepaspectratio]{#1}%
    }{%
        \fbox{%
            \begin{minipage}[c][#2][c]{0.96\linewidth}
                \centering
                Ảnh giao diện sẽ được bổ sung tại:\\[0.3em]
                \texttt{\detokenize{#1}}
            \end{minipage}%
        }%
    }%
}

\newcommand{\frontendpair}[6]{%
    \begin{figure}[H]
        \centering
        \begin{subfigure}[t]{0.48\textwidth}
            \centering
            \frontendimage{#1}{0.24\textheight}
            \caption{#2}
            \label{#3}
        \end{subfigure}\hfill
        \begin{subfigure}[t]{0.48\textwidth}
            \centering
            \frontendimage{#4}{0.24\textheight}
            \caption{#5}
            \label{#6}
        \end{subfigure}
    \end{figure}%
}

\newcommand{\frontendtriple}[9]{%
    \begin{figure}[H]
        \centering
        \begin{subfigure}[t]{0.31\textwidth}
            \centering
            \frontendimage{#1}{0.18\textheight}
            \caption{#2}
            \label{#3}
        \end{subfigure}\hfill
        \begin{subfigure}[t]{0.31\textwidth}
            \centering
            \frontendimage{#4}{0.18\textheight}
            \caption{#5}
            \label{#6}
        \end{subfigure}\hfill
        \begin{subfigure}[t]{0.31\textwidth}
            \centering
            \frontendimage{#7}{0.18\textheight}
            \caption{#8}
            \label{#9}
        \end{subfigure}
    \end{figure}%
}

\newcommand{\frontendtwoone}[9]{%
    \begin{figure}[H]
        \centering
        \begin{subfigure}[t]{0.48\textwidth}
            \centering
            \frontendimage{#1}{0.22\textheight}
            \caption{#2}
            \label{#3}
        \end{subfigure}\hfill
        \begin{subfigure}[t]{0.48\textwidth}
            \centering
            \frontendimage{#4}{0.22\textheight}
            \caption{#5}
            \label{#6}
        \end{subfigure}

        \par\medskip
        \begin{subfigure}[t]{0.48\textwidth}
            \centering
            \frontendimage{#7}{0.22\textheight}
            \caption{#8}
            \label{#9}
        \end{subfigure}
    \end{figure}%
}

\newcommand{\frontendonetwo}[9]{%
    \begin{figure}[H]
        \centering
        \begin{subfigure}[t]{0.48\textwidth}
            \centering
            \frontendimage{#1}{0.22\textheight}
            \caption{#2}
            \label{#3}
        \end{subfigure}

        \par\medskip
        \begin{subfigure}[t]{0.48\textwidth}
            \centering
            \frontendimage{#4}{0.22\textheight}
            \caption{#5}
            \label{#6}
        \end{subfigure}\hfill
        \begin{subfigure}[t]{0.48\textwidth}
            \centering
            \frontendimage{#7}{0.22\textheight}
            \caption{#8}
            \label{#9}
        \end{subfigure}
    \end{figure}%
}

\newcommand{\frontendsingle}[3]{%
    \begin{figure}[H]
        \centering
        \begin{subfigure}[t]{0.62\textwidth}
            \centering
            \frontendimage{#1}{0.25\textheight}
            \caption{#2}
            \label{#3}
        \end{subfigure}
    \end{figure}%
}

\paragraph{Tổng quan hiện thực Frontend.}

Frontend được hiện thực bằng React và tổ chức theo bốn lớp trách nhiệm: định tuyến, trang chức năng,
component dùng lại và lớp API. Các lớp này phối hợp để hiển thị dữ liệu, tiếp nhận thao tác và gửi
request đến Backend; quyền truy cập và dữ liệu nghiệp vụ vẫn được xác minh ở phía máy chủ. Khu vực
quản trị sử dụng tuyến riêng, còn các chức năng khám phá sản phẩm, tư vấn AI và đọc bài viết được
đặt trong các tuyến công khai hoặc yêu cầu đăng nhập tùy theo nghiệp vụ.

\begin{table}[H]
    \centering
    \caption{Ánh xạ các nhóm giao diện Frontend với chức năng chính}
    \label{tab:frontend-interface-map}
    \begin{tabularx}{\textwidth}{|>{\RaggedRight\arraybackslash}p{0.22\textwidth}|>{\RaggedRight\arraybackslash}p{0.25\textwidth}|X|}
        \hline
        \textbf{Nhóm giao diện} & \textbf{Tuyến tiêu biểu} & \textbf{Chức năng hiện thực} \\
        \hline
        Khám phá sản phẩm & \texttt{/}, \texttt{/category} & Hiển thị danh mục, danh sách sản phẩm, tìm kiếm, lọc và điều hướng đến chi tiết sản phẩm. \\
        \hline
        Hỗ trợ lựa chọn & \texttt{/prodetail/:id}, \texttt{/compare}, \texttt{/wishlist} & Xem chi tiết, so sánh, lưu yêu thích và xem nhanh sản phẩm. \\
        \hline
        Trợ lý AI & \texttt{/tu-van}, \texttt{/ai} & Hội thoại nhiều lượt, gửi văn bản hoặc hình ảnh và hiển thị sản phẩm được đề xuất. \\
        \hline
        Mua hàng & \texttt{/cart}, \texttt{/checkout} & Quản lý giỏ hàng, nhập thông tin nhận hàng, chọn thanh toán và xác nhận đơn. \\
        \hline
        Tài khoản & \texttt{/login}, \texttt{/register}, \texttt{/info}, \texttt{/order-history} & Xác thực, cập nhật hồ sơ và theo dõi lịch sử đơn hàng. \\
        \hline
        Quản trị & \texttt{/admin} & Quản lý sản phẩm, khách hàng, đơn hàng, khuyến mãi, coupon và bài viết. \\
        \hline
    \end{tabularx}
\end{table}

\subsubsection{Giao diện dành cho khách hàng}

\paragraph{Trang chủ và danh mục sản phẩm.}

Trang chủ là điểm bắt đầu của hành trình mua sắm, kết hợp thanh điều hướng, ô tìm kiếm, banner,
danh mục và các sản phẩm nổi bật. Từ trang chủ, người dùng có thể chuyển đến danh mục hoặc nhập
truy vấn để tìm sản phẩm. Trang danh mục hiển thị danh sách theo nhóm, hỗ trợ lọc theo thuộc tính,
sắp xếp và phân trang; các thao tác này được giữ trong cùng một luồng để người dùng không mất ngữ
cảnh duyệt sản phẩm.

\frontendpair
{Images/Frontend/giao diện trang chủ của hệ thống.png}{Trang chủ của hệ thống}{fig:frontend-home}
{Images/Frontend/giao diện danh mục, tìm kiếm và lọc sản phẩm.png}{Danh mục, tìm kiếm và lọc sản phẩm}{fig:frontend-category}

\paragraph{Chi tiết sản phẩm và hỗ trợ lựa chọn.}

Trang chi tiết tập trung các thông tin cần thiết trước khi mua gồm hình ảnh, tên, giá, biến thể,
tồn kho, mô tả và đánh giá. Người dùng có thể chọn biến thể, theo dõi trạng thái còn hàng, thêm sản
phẩm vào giỏ hoặc chuyển sang danh sách yêu thích. Chức năng xem nhanh và so sánh giúp đối chiếu
nhiều lựa chọn mà không phải rời khỏi ngữ cảnh danh mục.

\frontendpair
{Images/Frontend/giao diện chi tiết sản phẩm.png}{Chi tiết sản phẩm}{fig:frontend-product-detail}
{Images/Frontend/giao diện trang sản phẩm yêu thích.png}{Danh sách sản phẩm yêu thích}{fig:frontend-wishlist}

\paragraph{Tìm kiếm và tư vấn sản phẩm bằng AI.}

Chatbot nổi hỗ trợ đặt câu hỏi nhanh trong quá trình duyệt, còn trang trợ lý chuyên biệt cung cấp
không gian cho hội thoại nhiều lượt. Người dùng có thể gửi câu hỏi bằng văn bản hoặc ảnh tham chiếu;
Frontend hiển thị trạng thái gửi, đang xử lý, kết quả và cảnh báo khi dịch vụ AI không sẵn sàng.
Các sản phẩm được đề xuất nằm trong ngữ cảnh câu trả lời, cho phép mở trang chi tiết và tiếp tục
mua sắm. Giá và tồn kho hiển thị trên giao diện vẫn được lấy lại từ Backend.

\frontendpair
{Images/Frontend/giao diện chatbot.png}{Chatbot hỗ trợ mua sắm}{fig:frontend-chatbot}
{Images/Frontend/giao diện trợ lý mua sắm đa phương tiện.png}{Trợ lý mua sắm đa phương thức}{fig:frontend-ai-consulting}

\paragraph{Giỏ hàng và thanh toán.}

Giỏ hàng cho phép cập nhật số lượng, xóa mặt hàng và kiểm tra tổng tiền trước khi checkout. Tại
màn hình thanh toán, người dùng nhập thông tin nhận hàng, áp dụng mã giảm giá hợp lệ và kiểm tra
lại danh sách sản phẩm trước khi xác nhận. Sau khi Backend ghi nhận đơn COD thành công, giao diện
hiển thị kết quả và cung cấp đường dẫn đến lịch sử đơn hàng.

\frontendpair
{Images/Frontend/giao diện quản lý giỏ hàng.png}{Quản lý giỏ hàng}{fig:frontend-cart}
{Images/Frontend/giao diện thanh toán.png}{Checkout và lựa chọn phương thức thanh toán}{fig:frontend-checkout}

\frontendpair
{Images/Frontend/giao diện đặt hàng thành công.png}{Xác nhận đặt hàng thành công}{fig:frontend-order-success}
{Images/Frontend/giao diện lịch sử đơn hàng.png}{Lịch sử và trạng thái đơn hàng}{fig:frontend-order-history}

\paragraph{Đăng nhập, đăng ký và quản lý tài khoản.}

Giao diện đăng nhập cho phép khách hàng truy cập các chức năng cá nhân hóa; giao diện đăng ký kiểm
tra dữ liệu đầu vào trước khi tạo tài khoản. Sau khi đăng nhập, người dùng có thể cập nhật thông tin
cá nhân, quản lý địa chỉ và tiếp tục các thao tác mua sắm theo quyền được cấp. Các trạng thái tải,
lỗi xác thực và lỗi kết nối được hiển thị ngay tại biểu mẫu để người dùng biết cách xử lý.

\frontendtwoone
{Images/Frontend/giao diện đăng nhập.png}{Đăng nhập hệ thống}{fig:frontend-login}
{Images/Frontend/giao diện đăng ký.png}{Đăng ký tài khoản}{fig:frontend-register}
{Images/Frontend/giao diện tài khoản.png}{Quản lý thông tin tài khoản}{fig:frontend-account}

\subsubsection{Giao diện dành cho quản trị viên}

Khu vực quản trị sử dụng bố cục riêng và chỉ cho phép tài khoản có vai trò phù hợp truy cập. Thanh
điều hướng bên giúp chuyển giữa các nhóm nghiệp vụ; các màn hình danh sách dùng bảng, bộ lọc và
phân trang, trong khi thao tác tạo, sửa hoặc chuyển trạng thái được thực hiện qua biểu mẫu và hộp
thoại xác nhận.

\paragraph{Tổng quan quản trị và quản lý dữ liệu.}

Trang tổng quan cung cấp các chỉ số vận hành để quản trị viên nắm nhanh tình trạng sản phẩm, người
dùng và đơn hàng. Màn hình quản lý sản phẩm hỗ trợ theo dõi catalog và cập nhật thông tin hàng hóa;
màn hình quản lý người dùng phục vụ kiểm soát tài khoản, vai trò và trạng thái hoạt động.

\frontendonetwo
{Images/Frontend/giao diện tổng quan dành cho quản trị viên.png}{Tổng quan dành cho quản trị viên}{fig:frontend-admin-dashboard}
{Images/Frontend/giao diện quản lý sản phẩm.png}{Quản lý sản phẩm}{fig:frontend-admin-products}
{Images/Frontend/giao diện quản lý người dùng.png}{Quản lý người dùng}{fig:frontend-admin-users}

\paragraph{Quản lý đơn hàng.}

Màn hình đơn hàng cho phép tìm kiếm, xem chi tiết và chuyển trạng thái theo vòng đời đã thiết kế.
Thông tin khách hàng, mặt hàng, tổng tiền và lịch sử xử lý được trình bày cùng ngữ cảnh để hạn chế
thao tác nhầm. Việc thay đổi trạng thái được gửi đến Backend, nơi tiếp tục kiểm tra quyền và điều
kiện nghiệp vụ trước khi lưu.

\frontendsingle
{Images/Frontend/giao diện quản lý đơn hàng.png}{Quản lý đơn hàng}{fig:frontend-admin-orders}

\paragraph{Quản lý khuyến mãi và mã giảm giá.}

Màn hình khuyến mãi dùng để cấu hình chính sách giảm tự động theo thời gian hiệu lực và phạm vi áp
dụng. Mã giảm giá được quản lý riêng vì có vòng đời, điều kiện sử dụng và cách nhập khác với chương
trình khuyến mãi. Việc tách hai màn hình giúp quản trị viên kiểm soát các chương trình đang hoạt động
và giảm nhầm lẫn khi nhiều chính sách được áp dụng đồng thời.

\frontendpair
{Images/Frontend/giao diện quản lý khuyến mãi.png}{Quản lý chương trình khuyến mãi}{fig:frontend-admin-promotions}
{Images/Frontend/giao diện quản lý mã giảm giá.png}{Quản lý mã giảm giá}{fig:frontend-admin-coupons}

\paragraph{Quản lý bài viết.}

Giao diện bài viết hỗ trợ quản trị nội dung giới thiệu, hướng dẫn và tin tức của cửa hàng. Nội dung
được tách khỏi catalog để vòng đời biên tập không phụ thuộc vào thao tác cập nhật sản phẩm; biểu mẫu
cho phép tạo, sửa, xem trước và thay đổi trạng thái hiển thị theo quyền quản trị.

\frontendsingle
{Images/Frontend/giao diện quản lý bài viết.png}{Quản lý bài viết}{fig:frontend-admin-blogs}

\paragraph{Nguyên tắc phản hồi và khả năng sử dụng.}

Các giao diện thống nhất về phản hồi thao tác, trạng thái đang tải, trạng thái rỗng và thông báo lỗi.
Nút gửi được khóa trong thời gian xử lý để hạn chế request trùng; danh sách dài sử dụng phân trang
hoặc bộ lọc; thao tác ảnh hưởng đến đơn hàng và dữ liệu quản trị cần bước xác nhận. Việc tổ chức theo
component và kiểm tra trên nhiều kích thước màn hình giúp các biểu mẫu, bảng và hộp thoại vẫn dễ sử dụng.

% Đồng bộ bộ đếm sau bảng ltablex trong tệp giao diện; môi trường
% breakabletable hoàn nguyên một lần tăng thử của ltablex khi kết thúc.
\stepcounter{table}

\subsection{Chuẩn bị dữ liệu và hiện thực các thành phần AI}

Dữ liệu sản phẩm được thu thập tự động từ website Zara, gồm tên, mô tả, màu sắc, chất liệu, hình ảnh, danh mục và các thuộc tính liên quan. Do phần lớn dữ liệu ở dạng tiếng Anh, mô hình Qwen3.5-9B được dùng để Việt hóa các trường văn bản; những giá trị tiếng Anh cần thiết vẫn được giữ lại để bảo toàn thông tin nguồn. Sau khi rà soát và chuẩn hóa, 1.000 sản phẩm được chọn làm dữ liệu chính thức của đồ án và lưu trong tệp \texttt{1000\_final.sql}. Tệp này được dùng cho các chức năng sản phẩm, tìm kiếm, lọc và đồng bộ sang AI Service. Các lệnh chuẩn bị dữ liệu và huấn luyện được trình bày trong Phụ lục~\ref{appendix:source-code}.

\begin{breakabletable}
\centering
\caption{Quy trình chuẩn bị dữ liệu đầu vào cho database}
\label{tab:database-input-dataset}
\begin{tabularx}{\linewidth}{|p{0.30\linewidth}|X|}
\hline
\textbf{Thành phần} & \textbf{Mô tả} \\ \hline
Nguồn dữ liệu & Thu thập sản phẩm tự động từ website Zara \\ \hline
Việt hóa & Dịch trường văn bản bằng Qwen3.5-9B, giữ lại giá trị tiếng Anh cần thiết \\ \hline
Lựa chọn & Chọn 1.000 sản phẩm cho bộ dữ liệu của đồ án \\ \hline
Hoàn thiện & Rà soát, chuẩn hóa thuộc tính và cập nhật database \\ \hline
Kết quả & Tệp \texttt{1000\_final.sql} được dùng làm database chính thức \\ \hline
\end{tabularx}
\end{breakabletable}

\paragraph{Fine-tune và đóng gói mô hình YOLOv26.}
Để thích nghi với dữ liệu thương mại điện tử thời trang, YOLOv26 được fine-tune trên tập dữ liệu xây dựng từ Fashionpedia, một bộ dữ liệu ảnh thời trang có annotation cho bài toán phát hiện đối tượng \parencite{jia2020fashionpedia}. Dữ liệu tham khảo từ phiên bản trên Hugging Face Datasets, được lưu ở dạng Parquet và có giấy phép CC-BY-4.0 \parencite{huggingfaceFashionpedia}. Sau khi chuyển sang định dạng YOLO, các nhãn được gom thành bốn lớp: \textit{accessories}, \textit{bags}, \textit{clothing} và \textit{shoes}; đây là các nhóm phù hợp với sản phẩm trong hệ thống và đủ tổng quát cho tìm kiếm bằng hình ảnh.

\iffalse
\begin{lstlisting}[style=terminal,caption={Cài đặt thư viện Ultralytics phục vụ fine-tune YOLOv26},label={lst:install-ultralytics}]
pip install -U ultralytics
\end{lstlisting}
\fi

Tập dữ liệu gồm ba tập train, validation và test. File cấu hình xác định đường dẫn, bốn lớp nói trên và được dùng cùng thư viện \texttt{ultralytics} để huấn luyện.

\iffalse
\begin{lstlisting}[language=Python,caption={Tạo file cấu hình dữ liệu cho fine-tune YOLOv26},label={lst:yolo-data-yaml}]
import os

yaml_text = """
train: /kaggle/input/datasets/anhdngltrn/dungne/fashion_dataset/train/images
val: /kaggle/input/datasets/anhdngltrn/dungne/fashion_dataset/val/images
test: /kaggle/input/datasets/anhdngltrn/dungne/fashion_dataset/test/images

nc: 4
names:
  0: accessories
  1: bags
  2: clothing
  3: shoes
""".strip()

with open("/kaggle/working/data.yaml", "w") as f:
    f.write(yaml_text)

print("Saved /kaggle/working/data.yaml")
print(open("/kaggle/working/data.yaml").read())
\end{lstlisting}
\fi

Mô hình YOLOv26s được nạp từ weight \texttt{yolo26s.pt} và huấn luyện trong 40 epoch với ảnh đầu vào \texttt{640}, kế thừa weight đã học trước, dừng sớm sau 8 epoch không cải thiện và dùng mixed precision trên GPU.

\iffalse
\begin{lstlisting}[language=Python,caption={Fine-tune mô hình YOLOv26s trên tập dữ liệu thời trang},label={lst:yolo-train}]
import torch
from ultralytics import YOLO

DATA_YAML = "/kaggle/working/data.yaml"

print("CUDA available:", torch.cuda.is_available())
print("GPU count:", torch.cuda.device_count())

for i in range(torch.cuda.device_count()):
    print(f"GPU {i}: {torch.cuda.get_device_name(i)}")

ngpu = torch.cuda.device_count()
device_cfg = [0, 1] if ngpu >= 2 else 0
batch_size = 24
worker_size = 12

model = YOLO("yolo26s.pt")

results = model.train(
    data=DATA_YAML,
    epochs=40,
    imgsz=640,
    batch=batch_size,
    device=device_cfg,
    workers=worker_size,
    pretrained=True,
    patience=8,
    project="/kaggle/working",
    name="yolo26s_run2",
    val=True,
    amp=True,
    save=True,
    save_period=1,
    verbose=True
)
\end{lstlisting}
\fi

Kết quả được đánh giá bằng Precision, Recall, mAP50 và mAP50-95 như trình bày trong \Cref{tab:yolov26-finetune-result}.

\begin{breakabletable}
\centering
\caption{Kết quả đánh giá mô hình YOLOv26 sau khi fine-tune trên tập dữ liệu thời trang}
\label{tab:yolov26-finetune-result}
\begin{tabularx}{\linewidth}{|>{\RaggedRight\arraybackslash}p{0.13\linewidth}|>{\centering\arraybackslash}p{0.08\linewidth}|>{\centering\arraybackslash}p{0.14\linewidth}|>{\centering\arraybackslash}p{0.13\linewidth}|>{\centering\arraybackslash}p{0.10\linewidth}|>{\centering\arraybackslash}p{0.10\linewidth}|>{\centering\arraybackslash}X|}
\hline
\textbf{Lớp} & \textbf{Ảnh} & \textbf{Đối tượng} & \textbf{Precision} & \textbf{Recall} & \textbf{mAP50} & \textbf{mAP50--95} \\ \hline
Tất cả & 2339 & 8226 & 0.875 & 0.760 & 0.849 & 0.685 \\ \hline
Accessories & 932 & 1428 & 0.806 & 0.587 & 0.695 & 0.512 \\ \hline
Bags & 367 & 384 & 0.842 & 0.694 & 0.807 & 0.604 \\ \hline
Clothing & 2339 & 3966 & 0.926 & 0.882 & 0.947 & 0.861 \\ \hline
Shoes & 1267 & 2448 & 0.924 & 0.880 & 0.945 & 0.763 \\ \hline
\end{tabularx}
\end{breakabletable}

Mô hình đạt Precision 0.875, Recall 0.760, mAP50 0.849 và mAP50-95 0.685. Hai lớp \textit{clothing} và \textit{shoes} tốt nhất theo mAP50 (lần lượt 0.947 và 0.945), trong khi \textit{accessories} thấp hơn với mAP50-95 là 0.512 do vật thể nhỏ và dễ lẫn với nền. Tuy vậy, kết quả tổng thể đáp ứng mục tiêu phát hiện vùng sản phẩm để crop trước khi tạo image embedding.

Thời gian suy luận trung bình khoảng 1.8ms mỗi ảnh. Weight tốt nhất được lưu trong thư mục model của AI Service và nạp khi container khởi động.

\iffalse
\begin{lstlisting}[style=terminal,caption={Đường dẫn lưu weight YOLOv26 sau khi fine-tune},label={lst:yolo-weight-path}]
service-ai/Resource/best.pt
\end{lstlisting}
\fi

Khi khởi động, AI Service load weight này tại đường dẫn container \path{/app/Resource/best.pt}. Trong tìm kiếm bằng hình ảnh, service nhận ảnh từ Spring Boot Backend, phát hiện và crop vùng sản phẩm, tạo image embedding rồi truy vấn Qdrant. Nếu không phát hiện được vùng phù hợp hoặc độ tin cậy thấp, ảnh gốc được dùng làm fallback để không gián đoạn quá trình tạo embedding.
\subsection{Cấu hình môi trường và chạy hệ thống local}

Hệ thống hỗ trợ hai chế độ vận hành: chạy toàn bộ trên máy local để phát triển, kiểm thử và triển khai phân tán trên Google Cloud Platform (GCP) để phục vụ bản demo. Cả hai chế độ đều sử dụng Docker để bảo đảm môi trường thực thi nhất quán, tuy nhiên vị trí của AI Service và cách kết nối giữa các dịch vụ có sự khác biệt.

Các lệnh chuẩn bị biến môi trường, khởi động Docker Compose và kiểm tra nhanh được đưa vào Phụ
lục~\ref{appendix:source-code}; phần này tập trung giải thích vai trò của từng dịch vụ và thứ tự khởi động.

\begin{breakabletable}
\centering
\caption{So sánh hai chế độ vận hành của hệ thống}
\label{tab:deployment-modes}
\begin{tabularx}{\linewidth}{|p{0.20\linewidth}|X|X|}
\hline
\textbf{Nội dung} & \textbf{Chạy local} & \textbf{Triển khai trên GCP} \\ \hline
Thành phần thực thi & MySQL, Qdrant, AI Service, Spring Boot Backend và Frontend cùng chạy bằng Docker Compose trên một máy & Frontend, Spring Boot Backend, MySQL và Qdrant chạy trên GCE VM; AI Service chạy riêng trên Cloud Run \\ \hline
Kết nối dịch vụ & Các container giao tiếp qua mạng mặc định của Docker Compose bằng tên dịch vụ & Cloud Run kết nối đến Backend và Qdrant trên địa chỉ private của VM thông qua VPC Connector \\ \hline
Mô hình AI & Bedrock cung cấp LLM và Cohere Embed v4 với vector văn bản, hình ảnh 1.536 chiều & GPU Cloud Run tự phục vụ Qwen3-Embedding-8B 4.096 chiều, Qwen3-VL-Embedding-2B 2.048 chiều và YOLO; Bedrock chỉ cung cấp LLM \\ \hline
Mục đích & Phát triển, kiểm thử chức năng và đánh giá mô hình & Vận hành bản demo và kiểm thử hệ thống trong môi trường đám mây \\ \hline
Dữ liệu lâu dài & Docker named volume cho MySQL và Qdrant & Docker named volume trên VM; cấu hình và khóa bí mật được tách khỏi mã nguồn \\ \hline
\end{tabularx}
\end{breakabletable}

Các công nghệ và phiên bản nền tảng chính được sử dụng gồm Java 21 và Spring Boot 4 cho Backend, Python 3.11 cùng FastAPI cho AI Service, React 19 và Vite cho Frontend, MySQL 8.4, Qdrant và Docker Compose. Frontend được build bằng Node.js 20 rồi phục vụ dưới dạng tệp tĩnh qua nginx. Backend và Frontend sử dụng Dockerfile nhiều giai đoạn nhằm giảm kích thước image chạy thực tế.

\paragraph{Đóng gói và chạy các dịch vụ tại local.}

Ở chế độ local, tệp \texttt{docker-compose.yml} tại thư mục gốc mô tả toàn bộ các dịch vụ. MySQL và Qdrant sử dụng named volume lần lượt là \texttt{mysql\_data} và \texttt{qdrant\_storage}. Chỉ Frontend tại cổng 80 và Spring Boot Backend tại cổng 8080 được ánh xạ ra máy host; MySQL, Qdrant và AI Service chỉ giao tiếp trong mạng Docker để hạn chế truy cập trực tiếp không cần thiết.

Trước khi khởi động, hai tệp cấu hình mẫu của Backend và AI Service được sao chép thành tệp \texttt{.env}. Người triển khai cần cập nhật mật khẩu MySQL, JWT secret, internal API key và thông tin xác thực của các dịch vụ mô hình. Giá trị \texttt{AI\_SERVICE\_INTERNAL\_API\_KEY} ở Backend phải trùng với \texttt{INTERNAL\_API\_KEY} và \texttt{SPRINGBOOT\_API\_KEY} ở AI Service. Các khóa bí mật chỉ được đặt trong biến môi trường hoặc tệp secret cục bộ và không được đưa vào Git.

\iffalse
\begin{lstlisting}[style=terminal,caption={Chuẩn bị cấu hình để chạy hệ thống tại local},label={lst:local-env-prepare}]
Copy-Item .\service-api-ecommerce\.env.example .\service-api-ecommerce\.env
Copy-Item .\service-ai\.env.example .\service-ai\.env

$env:MYSQL_ROOT_PASSWORD = 'local-mysql-password'
$env:BEDROCK_REGION = 'ap-northeast-1'
$env:AWS_BEARER_TOKEN_BEDROCK = 'bedrock-api-key'

docker compose config --quiet
\end{lstlisting}
\fi

Với volume MySQL mới, tệp \texttt{1000\_final.sql} được import một lần để khởi tạo dữ liệu chính thức của đồ án. Dữ liệu không được import lại ở mỗi lần khởi động vì có thể ghi đè các thay đổi phát sinh trong quá trình sử dụng.

\iffalse
\begin{lstlisting}[style=terminal,caption={Khởi động các dịch vụ và khởi tạo dữ liệu local},label={lst:local-compose-start}]
docker compose up -d mysql qdrant

Get-Content -Raw .\database\1000_final.sql |
  docker compose exec -T mysql sh -c 'exec mysql -uroot --password=$MYSQL_ROOT_PASSWORD DACN'

docker compose up -d --build service-ai
docker compose up -d --build spring-api frontend
docker compose ps
\end{lstlisting}
\fi

Thứ tự trên bảo đảm hai dịch vụ dữ liệu sẵn sàng trước AI Service và Backend. AI Service khởi tạo các collection còn thiếu trong Qdrant khi container bắt đầu chạy. Sau đó, chức năng đồng bộ AI của Backend đọc dữ liệu sản phẩm từ MySQL, tạo embedding và cập nhật các vector tương ứng vào Qdrant. Việc đồng bộ có thể áp dụng cho toàn bộ 1.000 sản phẩm ở local hoặc giới hạn số lượng khi chỉ cần kiểm thử nhanh.

Khi hệ thống đã sẵn sàng, Frontend được truy cập tại \texttt{http://localhost}, còn Backend tại \texttt{http://localhost:8080}. Trạng thái AI Service được xác nhận qua health status của container vì cổng 8000 không được publish ra host trong cấu hình Compose gốc.

\iffalse
\begin{lstlisting}[style=terminal,caption={Kiểm tra hệ thống chạy local},label={lst:local-compose-check}]
docker compose ps
Invoke-WebRequest http://localhost:8080/api/blogs
Invoke-WebRequest http://localhost
\end{lstlisting}
\fi

\subsection{Triển khai hệ thống trên Google Cloud Platform}

Bản GCP được triển khai theo mô hình phân tán: GCE VM \texttt{datn-stack-vm} tại zone \texttt{asia-southeast1-b} chạy Frontend, Spring Boot Backend, MySQL và Qdrant; AI Service chạy độc lập trên Cloud Run tại region \texttt{asia-southeast1}. Cách tổ chức này cho phép cập nhật và điều chỉnh tài nguyên GPU mà không ảnh hưởng cụm nghiệp vụ. Các câu lệnh build, deploy và kiểm tra được chuyển sang Phụ lục~\ref{appendix:source-code}; phần này chỉ trình bày cấu hình đã nghiệm thu.

\clearpage
\begin{breakabletable}
\centering
\caption{Cấu hình hạ tầng triển khai trên GCP}
\label{tab:gcp-resource-config}
\begin{tabularx}{\linewidth}{|p{0.22\linewidth}|p{0.25\linewidth}|X|}
\hline
\textbf{Tài nguyên} & \textbf{Cấu hình} & \textbf{Vai trò} \\ \hline
GCE VM & \texttt{e2-standard-4}, Ubuntu 22.04, đĩa 150 GB & Chạy Frontend, Backend, MySQL và Qdrant bằng Docker Compose \\ \hline
Cloud Run GPU & NVIDIA RTX PRO 6000 Blackwell, 20 vCPU, 80 GiB RAM, concurrency 1, tối đa một instance & Chạy FastAPI, hai mô hình Qwen embedding và YOLO trên CUDA; gọi Bedrock cho tác vụ LLM \\ \hline
Artifact Registry & Docker repository tại \texttt{asia-southeast1} & Lưu trữ các image đã được gắn tag bất biến \\ \hline
VPC Connector & Dải mạng \texttt{10.8.0.0/28} & Cho phép Cloud Run truy cập Backend và Qdrant qua mạng private \\ \hline
Memorystore & Redis tại địa chỉ private trong VPC & Lưu sở thích người dùng lâu dài cho trợ lý hội thoại \\ \hline
Secret Manager & Secret được tách theo chức năng và cấp cho runtime service account & Quản lý Bedrock bearer token và internal API key, không ghi secret vào image hoặc repository \\ \hline
\end{tabularx}
\end{breakabletable}

Frontend và Backend lần lượt được công khai qua cổng 80 và 8080. MySQL không publish cổng 3306; Qdrant chỉ bind cổng 6333 trên private IP và firewall chỉ nhận lưu lượng từ VPC Connector. Cloud Run dùng ingress nội bộ để trao đổi với Backend và Qdrant mà không công khai tầng dữ liệu ra Internet. Bedrock bearer token và internal API key được ánh xạ từ Secret Manager vào runtime service account, không ghi trong image, câu lệnh deploy hoặc repository.

AI Service tự phục vụ embedding trên GPU: \texttt{Qwen/Qwen3-Embedding-8B} tạo vector văn bản 4.096 chiều và \texttt{Qwen/Qwen3-VL-Embedding-2B} tạo vector ảnh 2.048 chiều. Hai mô hình chạy CUDA với \texttt{bfloat16}; vector được chuẩn hóa và truy vấn bằng độ đo Cosine. Weight YOLO đã fine-tune tại \path{Resource/best.pt} phát hiện và crop vật thể trước khi tạo image embedding, đồng thời dùng ảnh gốc làm fallback. Amazon Bedrock không tạo embedding mà chỉ chạy \texttt{openai.gpt-oss-120b-1:0} tại \texttt{ap-northeast-1} để phân tích ý định, trích xuất điều kiện và sinh phản hồi. Reranker được tắt. Text Search kết hợp độ tương đồng vector với tín hiệu lexical và metadata; Image Search dùng độ tương đồng ảnh; truy vấn ảnh--văn bản dùng ảnh làm tín hiệu xếp hạng chính, áp dụng ràng buộc văn bản và chỉ truy hồi văn bản để phục hồi khi thiếu ứng viên, không fusion điểm text--image.

\paragraph{Triển khai các dịch vụ trên GCE VM.}

Mã nguồn được đặt tại \path{/opt/dATN}; biến production nằm trong \path{.env.production} với quyền truy cập hạn chế. Tệp \path{docker-compose.vm.yml} ghi đè cấu hình local để tắt AI container trên VM, bind Qdrant vào private IP và trỏ Backend tới Cloud Run. Khi cập nhật, cấu hình Compose được kiểm tra trước; MySQL và Qdrant khởi động trước, tiếp theo là Backend rồi Frontend. Hai named volume \texttt{datn\_mysql\_data} và \texttt{datn\_qdrant\_storage} được giữ độc lập với vòng đời container. Do \texttt{VITE\_API\_BASE\_URL} được đóng vào bundle, Frontend phải build lại nếu địa chỉ Backend thay đổi.

\paragraph{Triển khai AI Service trên Cloud Run.}

Cloud Build tạo image từ \path{Dockerfile.gpu} và \path{cloudbuild.gpu.kaniko.yaml} bằng Kaniko trên worker \texttt{E2\_HIGHCPU\_8}. Snapshot của hai mô hình Qwen được tải sẵn; image được gắn tag bất biến và lưu digest để truy vết, rollback. Revision nghiệm thu \path{service-ai-gptoss-pilot-qwen-bw-v2} dùng execution environment thế hệ hai, VPC Connector, một GPU NVIDIA RTX PRO 6000 Blackwell, 20 vCPU, 80 GiB RAM, concurrency 1, timeout 3.600 giây và tối đa một instance. Sau khi \path{/health/ready} xác nhận text embedding, image embedding, YOLO và LLM đều sẵn sàng, toàn bộ traffic được chuyển sang revision này.

Trong khi đồng bộ và demo, \path{min-instances=1} giữ model trong GPU; sau nghiệm thu giá trị được trả về 0 để scale-to-zero, còn \path{max-instances=1}. Memorystore lưu hồ sơ sở thích lâu dài, nhưng không có RedisJSON nên checkpoint LangGraph của bản demo được lưu bằng SQLite trên filesystem tạm. Vì checkpoint có thể mất khi instance dừng, cấu hình dài hạn cần chuyển sang PostgreSQL hoặc kho bền vững tương thích.

\subsection{Đồng bộ dữ liệu, kiểm tra và khôi phục}

MySQL là nguồn dữ liệu chuẩn; Qdrant chỉ giữ vector và metadata dẫn xuất. Bản GCP dùng collection văn bản \path{product_text_vectors_qwen3_embedding_8b_4096_metadata_v4_1000_products} và collection ảnh \path{product_image_vectors_qwen3_vl_embedding_2b_2048_1000_products}. Hai collection dùng Cosine và được tách khỏi collection Bedrock 1.536 chiều của môi trường local để không trộn vector khác mô hình hoặc số chiều.

Sau khi Backend và AI Service sẵn sàng, tác vụ full sync đọc 1.000 sản phẩm theo batch. Mỗi sản phẩm tạo một text vector; ảnh hợp lệ được YOLO phát hiện, crop rồi tạo vector theo vùng. Tác vụ hoàn thành sau khoảng 51 phút. Lần đầu có 999 sản phẩm thành công và một lỗi kết nối Qdrant tạm thời; hệ thống chỉ retry sản phẩm cùng các URL timeout, sau đó đạt kết quả ở Bảng~\ref{tab:gcp-sync-result}.

\begin{breakabletable}
\centering
\caption{Kết quả đồng bộ và đối soát dữ liệu trên GCP}
\label{tab:gcp-sync-result}
\begin{tabularx}{\linewidth}{|p{0.34\linewidth}|X|}
\hline
\textbf{Chỉ tiêu} & \textbf{Kết quả nghiệm thu} \\ \hline
Sản phẩm trong MySQL & 1.000 sản phẩm; cả 1.000 sản phẩm có dữ liệu ảnh \\ \hline
Text collection & Trạng thái \texttt{green}; 1.000 point; 1.000 product ID duy nhất; vector 4.096 chiều \\ \hline
Image collection & Trạng thái \texttt{green}; 13.879 point; 1.000 product ID duy nhất; vector 2.048 chiều \\ \hline
URL ảnh đã lập chỉ mục & 8.871 URL ảnh hợp lệ \\ \hline
Độ phủ sản phẩm & 1.000/1.000 sản phẩm có text vector và 1.000/1.000 sản phẩm có image vector \\ \hline
\end{tabularx}
\end{breakabletable}

Trong 8.872 URL media của MySQL có một playlist HLS \texttt{master.m3u8}, không phải ảnh nên bị loại; sản phẩm tương ứng vẫn có tám image vector. Vì vậy 8.871 URL hợp lệ vẫn bao phủ đủ 1.000 sản phẩm. Nghiệm thu đối soát độc lập số product ID, số point, kích thước vector và cặp \texttt{product\_id--image\_url}, thay vì chỉ dựa vào trạng thái tác vụ.

\paragraph{Kiểm tra sau cài đặt và triển khai.}

Hệ thống được kiểm tra từ tầng dữ liệu đến giao diện theo các tiêu chí sau:

\begin{itemize}
    \item MySQL, Qdrant và hai named volume hoạt động; Backend truy cập được MySQL và AI Service.
    \item \path{/health} và \path{/health/ready} hợp lệ; hai collection đúng kích thước vector và bao phủ đủ 1.000 product ID.
    \item Tìm kiếm text dùng Qwen3-Embedding-8B; tìm kiếm ảnh dùng YOLO và Qwen3-VL-Embedding-2B; hybrid trả tối đa năm sản phẩm.
    \item Trợ lý nhận phản hồi tiếng Việt từ Bedrock; Frontend gọi đúng Backend và không có lỗi xác thực secret.
\end{itemize}
Log VM được theo dõi bằng \texttt{docker logs}, log AI bằng Cloud Logging và không in toàn bộ environment để tránh lộ secret. Smoke test xác nhận text, image và hybrid dùng đúng collection Qwen; YOLO phát hiện hai vùng trên ảnh thử mà không cần fallback; trợ lý nhận phản hồi Bedrock. Chi phí build, giữ GPU để đồng bộ và kiểm thử ước tính 5,5--7 USD, chưa gồm phí Bedrock nhỏ; sau nghiệm thu GPU được scale-to-zero.

\paragraph{Sao lưu, rollback và khôi phục khi triển khai.}

Trước khi deploy hoặc đổi schema, MySQL được dump và Qdrant được snapshot; không dùng \texttt{docker compose down -v} hoặc xóa named volume. Nếu AI revision lỗi, Cloud Run chuyển traffic về revision tốt gần nhất; Backend và Frontend trở về image tag đã xác nhận. MySQL chỉ được thay sau khi bản phục hồi vượt kiểm tra trên database tạm, còn collection Qdrant cũ được giữ đến khi collection mới đồng bộ và vượt smoke test.

\subsection*{Kết chương}

Chương này đã trình bày quá trình chuyển thiết kế thành hệ thống có thể chạy được,
từ hiện thực giao diện và dữ liệu đến đóng gói các dịch vụ, chạy local và triển
khai trên GCP. Các thành phần nghiệp vụ được giữ độc lập với AI Service; dữ liệu
MySQL được bảo toàn qua các lần cập nhật; chỉ mục Qdrant có thể đồng bộ lại từ
dữ liệu chuẩn. Quy trình health check, sao lưu và rollback giúp hệ thống có thể
được kiểm tra, cập nhật và khôi phục có kiểm soát trước khi đánh giá ở chương
tiếp theo.
